
\chapter{Типичные ошибки и диагностика}
\label{chap:mistakes}

Ошибки CIAO могут возникать при компиляции модели,
исполнении переходов или подключении пользовательского
интерфейса.

При диагностике сначала обратите внимание на тип
исключения и его сообщение. Если ошибка произошла
во время исполнения, дополнительно изучите трассу
Runtime: она поможет определить, какая команда
и какой переход привели к проблеме.

\section{Распространённые ошибки}

\begin{table}[ht]
\centering
\small
\caption{Частые ошибки CIAO}
\begin{tabularx}{\textwidth}{@{}p{0.36\textwidth}X@{}}
\toprule
Ошибка & Причина и решение \\
\midrule

Некорректный endpoint &
Операция принадлежит другому классу автомата.
Убедитесь, что в паре указан нужный объект
и операция его класса. \\

Отсутствует маршрут required-операции &
Не объявлен \code{Connection} или внешний обработчик.
Добавьте необходимую связь. \\

\code{NoTransitionError} &
Для команды нет перехода из текущего состояния.
Проверьте состояние автомата и допустимые события. \\

\code{AmbiguousTransitionError} &
Несколько Guard вернули \code{True}.
Сделайте условия взаимоисключающими. \\

Ошибка сигнатуры обработчика &
Проверьте типы параметров, возвращаемое значение
и наличие обязательных \code{Emit}/\code{Query}
при их объявлении в декораторе. \\

\code{UIBindingConflictError} &
Повторная привязка к кнопке или внешней операции.
Удалите конфликтующую регистрацию. \\

\code{RuntimeReentryError} &
Повторный вызов \code{step()} или \code{run()}
во время исполнения. Используйте механизм
планирования шагов UI-адаптера. \\

\bottomrule
\end{tabularx}
\end{table}

Не все ошибки обнаруживаются при компиляции.
Например, совместимость сигнатур можно проверить
заранее, а результат Guard становится известен
только во время исполнения.

Если исключение возникает внутри Action, Runtime
не фиксирует новое состояние и не публикует
подготовленные команды. Однако изменения контекста,
выполненные до ошибки, сохраняются.

\section{Рекомендации по организации автоматов}

При разработке собственных автоматов придерживайтесь
нескольких правил:

\begin{itemize}
  \item Не храните \code{Runtime} или \code{Window}
    в контексте автомата. Контекст предназначен
    для данных самого объекта.

  \item Не вызывайте provided- и required-команды
    как обычные Python-методы. Используйте
    \code{Runtime.send} и \code{Emit}/\code{Emits}.

  \item Не изменяйте контекст внутри Guard или Query.
    Для изменений используйте Action и команды.

  \item Разделяйте действия над контекстом и
    отправку команд на отдельные Action.

  \item Не обращайтесь из Controller напрямую
    к Tkinter-виджетам. Передавайте данные через View.
\end{itemize}

Эти правила помогают сохранять поведение автоматов
явным и упрощают чтение их диаграмм.

\section{Проверка программы без UI}

Если приложение работает неправильно, полезно
сначала проверить его автоматную часть без Tkinter.

Для этого достаточно собрать \code{Program}, создать
\code{Runtime} и отправить команды непосредственно
через \code{send()}.

Например, для Array Sum:

\begin{lstlisting}[language=Python,numbers=none]
runtime = Runtime(sample.program)

runtime.send(
    sample.view,
    SumView.init_clicked,
    request,
)

runtime.run(max_steps=100)
assert runtime.pending_messages == 0
\end{lstlisting}

Здесь предполагается, что \code{sample} содержит
собранную программу, а \code{request} --- подготовленные
параметры запуска.

Так можно проверить выполнение переходов и обмен
командами между автоматами без графического интерфейса.
Помимо завершения очереди, проверяйте ожидаемые
состояния и результаты вычисления.

Если автоматная часть работает корректно, переходите
к проверке UI-привязок: правильности передаваемых
параметров, обработчиков кнопок и методов отображения.

Так проще отделить ошибки самой модели от ошибок
её взаимодействия с интерфейсом.
